\begin{corrige}{3-7}

Soit $\Omega = \{ 1, 2, 3, 4, 5, 6 \}$ l'univers associé à l'expérience aléatoire. On suppose que le dé n'est pas truqué et qu'il y a donc équiprobabilité.

Soit $D$ l'évènement ``obtenir un $2$'' et $Np$ l'évènement ``obtenir un nombre pair''. On cherche en fait la probabilité conditionnelle:
\[
 P ( D | Np ) = \frac{P(D \cap Np)}{P(Np)}.
\]
Comme $2$ est un nombre pair, nous avons en fait: $D \cap Np = D$ et donc $P(D \cap Np) = P(D) = \frac16$. La probabilité d'avoir un nombre pair est donnée par $P(Np) = \frac36 = \frac12$ (c'est celle de sortir un $2$, un $4$ ou un $6$).
Ceci donne au final:
\[
 P ( D | Np ) = \frac{P(D \cap Np)}{P(Np)} = \frac{\frac16}{\frac12} = \frac13.
\]


\end{corrige}
